% GENERATED FILE. Edit canonical lessons, then run npm run generate:books.
% canonical-source-hash: fnv1a64:e7878a172e0ed2f1
% canonical-lessons: JA-C367-hakkiri, JA-C367-tsumari, JA-C367-juubun, JA-C367-teido, JA-C367-wariai

\chapter{Being Clear}
\label{ch:ja-being-clear}
\hlchaptermodality{367}

\begin{chapteropening}
\textbf{By the end of this chapter:} \emph{I can say a thing clearly and sum up what it means.}

Complete the last lesson of chapter 367: I can say a thing clearly and sum up what it means.
\end{chapteropening}

\section[\texorpdfstring{\ja{はっきり} (hakkiri)}{はっきり (hakkiri)}]{\ja{はっきり} (hakkiri) — clearly; plainly}

\label{lesson:JA-C367-hakkiri}



\begin{quote}
Before the new one: say the Japanese for to put first, then the Japanese for to be impatient.
\end{quote}

\subsection*{You'll want to know: \ja{はっきり}}
\textbf{\ja{はっきり}} — \emph{hakkiri} — \textquotedblleft{}clearly; plainly\textquotedblright{}.

\begin{quote}
\textbf{\ja{しつもんが} \ja{あったら、はっきり} \ja{言}\ja{って} \ja{ください。}} — \emph{shitsumon ga attara, hakkiri itte kudasai} — \textquotedblleft{}If you have a question, please say it plainly.\textquotedblright{}
\end{quote}

\begin{cousinweb}
\textbf{\ja{はっきり}} (\emph{hakkiri}) is a native mimetic word, one of the Japanese words that picture a manner or a feeling in sound; JMdict marks it as mimetic. It is inherited, not borrowed, and its older origin is not settled (JMdict).
\end{cousinweb}

\begin{grammarlens}[title={What you've built}]
One more word for saying what happens if.
\end{grammarlens}

\subsection*{Your turn}
\begin{itemize}
\raggedright
  \item \emph{Say it:} \emph{hakkiri}
  \item \emph{Say it:} \emph{hakkiri}, once more
  \item \emph{Say it:} \emph{shitsumon ga attara, hakkiri itte kudasai} — If you have a question, please say it plainly.
  \item \emph{Recall:} say the Japanese for to put first, then the Japanese for to be impatient, then say \emph{hakkiri} again
\end{itemize}

\begin{tcolorbox}[breakable,colback=teal!4,colframe=teal!35!black,title={Before you move on}]
What does \ja{はっきり} mean? (\textquotedblleft{}Clearly; plainly\textquotedblright{}.) Say it once more.
\end{tcolorbox}
\section[\texorpdfstring{\ja{つまり} (tsumari)}{つまり (tsumari)}]{\ja{つまり} (tsumari) — in other words; that is}

\label{lesson:JA-C367-tsumari}



\begin{quote}
Before the new one: say the Japanese for to be impatient, then the Japanese for clearly.
\end{quote}

\subsection*{You'll want to know: \ja{つまり}}
\textbf{\ja{つまり}} — \emph{tsumari} — \textquotedblleft{}in other words; that is\textquotedblright{}.

\begin{quote}
\textbf{\ja{つまり、あめが} \ja{ふったら、ちゅうしです。}} — \emph{tsumari, ame ga futtara, chūshi desu} — \textquotedblleft{}In other words, if it rains, it is called off.\textquotedblright{}
\end{quote}

\begin{cousinweb}
\textbf{\ja{つまり}} (\emph{tsumari}) is a native Japanese word (wago), inherited rather than borrowed from Chinese, and its older origin is not settled. JMdict notes that it is usually written in kana (JMdict).
\end{cousinweb}

\begin{grammarlens}[title={What you've built}]
One more word for saying what happens if.
\end{grammarlens}

\subsection*{Your turn}
\begin{itemize}
\raggedright
  \item \emph{Say it:} \emph{tsumari}
  \item \emph{Say it:} \emph{tsumari}, once more
  \item \emph{Say it:} \emph{tsumari, ame ga futtara, chūshi desu} — In other words, if it rains, it is called off.
  \item \emph{Recall:} say the Japanese for to be impatient, then the Japanese for clearly, then say \emph{tsumari} again
\end{itemize}

\begin{tcolorbox}[breakable,colback=teal!4,colframe=teal!35!black,title={Before you move on}]
What does \ja{つまり} mean? (\textquotedblleft{}In other words; that is\textquotedblright{}.) Say it once more.
\end{tcolorbox}
\section[\texorpdfstring{\ja{じゅうぶん} (jūbun)}{じゅうぶん (jūbun)}]{\ja{じゅうぶん} (jūbun) — enough; plenty}

\label{lesson:JA-C367-juubun}



\begin{quote}
Before the new one: say the Japanese for clearly, then the Japanese for in other words.
\end{quote}

\subsection*{You'll want to know: \ja{じゅうぶん}}
\textbf{\ja{じゅうぶん}} — \emph{jūbun} — \textquotedblleft{}enough; plenty\textquotedblright{}.

\begin{quote}
\textbf{\ja{じかんが} \ja{じゅうぶん} \ja{あったら、えきまで} \ja{あるきます。}} — \emph{jikan ga jūbun attara, eki made arukimasu} — \textquotedblleft{}If there is plenty of time, I walk to the station.\textquotedblright{}
\end{quote}

A na-adjective: \ja{な} before a noun.

\begin{cousinweb}
\textbf{\ja{じゅうぶん}} (\emph{jūbun}) is Sino-Japanese (kango): two Chinese characters read with their on'yomi, readings borrowed from Chinese. They are \emph{jū}, \textquotedblleft{}ten\textquotedblright{} (Middle Chinese \emph{dzyip}) and \emph{bun}, \textquotedblleft{}part\textquotedblright{} (\emph{pjun}) (KANJIDIC2; Middle Chinese in Baxter's transcription).
\end{cousinweb}

\begin{grammarlens}[title={What you've built}]
One more word for saying what happens if.
\end{grammarlens}

\subsection*{Your turn}
\begin{itemize}
\raggedright
  \item \emph{Say it:} \emph{jūbun}
  \item \emph{Say it:} \emph{jūbun}, once more
  \item \emph{Say it:} \emph{jikan ga jūbun attara, eki made arukimasu} — If there is plenty of time, I walk to the station.
  \item \emph{Recall:} say the Japanese for clearly, then the Japanese for in other words, then say \emph{jūbun} again
\end{itemize}

\begin{tcolorbox}[breakable,colback=teal!4,colframe=teal!35!black,title={Before you move on}]
What does \ja{じゅうぶん} mean? (\textquotedblleft{}Enough; plenty\textquotedblright{}.) Say it once more.
\end{tcolorbox}
\section[\texorpdfstring{\ja{ていど} (teido)}{ていど (teido)}]{\ja{ていど} (teido) — a degree; an extent}

\label{lesson:JA-C367-teido}



\begin{quote}
Before the new one: say the Japanese for in other words, then the Japanese for enough.
\end{quote}

\subsection*{You'll want to know: \ja{ていど}}
\textbf{\ja{ていど}} — \emph{teido} — \textquotedblleft{}a degree; an extent\textquotedblright{}.

\begin{quote}
\textbf{\ja{くすりは} \ja{ねつの} \ja{ていどに} \ja{よります。}} — \emph{kusuri wa netsu no teido ni yorimasu} — \textquotedblleft{}The medicine depends on how high the fever is.\textquotedblright{}
\end{quote}

\begin{cousinweb}
\textbf{\ja{ていど}} (\emph{teido}) is Sino-Japanese (kango): two Chinese characters read with their on'yomi, readings borrowed from Chinese. They are \emph{tei}, \textquotedblleft{}extent\textquotedblright{} (Middle Chinese \emph{drjeng}) and \emph{do}, \textquotedblleft{}degrees\textquotedblright{} (\emph{duH}) (KANJIDIC2; Middle Chinese in Baxter's transcription).
\end{cousinweb}

\begin{grammarlens}[title={What you've built}]
One more word for saying what happens if.
\end{grammarlens}

\subsection*{Your turn}
\begin{itemize}
\raggedright
  \item \emph{Say it:} \emph{teido}
  \item \emph{Say it:} \emph{teido}, once more
  \item \emph{Say it:} \emph{kusuri wa netsu no teido ni yorimasu} — The medicine depends on how high the fever is.
  \item \emph{Recall:} say the Japanese for in other words, then the Japanese for enough, then say \emph{teido} again
\end{itemize}

\begin{tcolorbox}[breakable,colback=teal!4,colframe=teal!35!black,title={Before you move on}]
What does \ja{ていど} mean? (\textquotedblleft{}A degree; an extent\textquotedblright{}.) Say it once more.
\end{tcolorbox}
\section[\texorpdfstring{\ja{わりあい} (wariai)}{わりあい (wariai)}]{\ja{わりあい} (wariai) — a proportion; a rate}

\label{lesson:JA-C367-wariai}



\begin{quote}
Before the new one: say the Japanese for enough, then the Japanese for a degree.
\end{quote}

\subsection*{You'll want to know: \ja{わりあい}}
\textbf{\ja{わりあい}} — \emph{wariai} — \textquotedblleft{}a proportion; a rate\textquotedblright{}.

\begin{quote}
\textbf{\ja{こうれいしゃの} \ja{わりあいが} \ja{ふえたら、ねんきんが} \ja{すくなく} \ja{なるでしょう。}} — \emph{kōreisha no wariai ga fuetara, nenkin ga sukunaku naru deshō} — \textquotedblleft{}If the share of elderly people grows, pensions will probably shrink.\textquotedblright{}
\end{quote}

\begin{cousinweb}
\textbf{\ja{わりあい}} (\emph{wariai}) is a native Japanese word (wago), inherited rather than borrowed from Chinese: it is built from \emph{wari}, \textquotedblleft{}dividing (waru)\textquotedblright{} and \emph{ai}, \textquotedblleft{}meeting (au)\textquotedblright{}, and its kanji were attached later as kun'yomi (JMdict).
\end{cousinweb}

\begin{grammarlens}[title={What you've built}]
One more word for saying what happens if.
\end{grammarlens}

\subsection*{Your turn}
\begin{itemize}
\raggedright
  \item \emph{Say it:} \emph{wariai}
  \item \emph{Say it:} \emph{wariai}, once more
  \item \emph{Say it:} \emph{kōreisha no wariai ga fuetara, nenkin ga sukunaku naru deshō} — If the share of elderly people grows, pensions will probably shrink.
  \item \emph{Recall:} say the Japanese for enough, then the Japanese for a degree, then say \emph{wariai} again
\end{itemize}

\begin{tcolorbox}[breakable,colback=teal!4,colframe=teal!35!black,title={Before you move on}]
What does \ja{わりあい} mean? (\textquotedblleft{}A proportion; a rate\textquotedblright{}.) Say it once more.
\end{tcolorbox}
